%%%%%%%%%%%%%%%% Header %%%%%%%%%%%%%%%%
\noindent\begin{minipage}{0.98\textwidth}
  \vskip 0mm
  \noindent
  { \begin{tabular}{p{7.5cm}}
      {\bfseries \sffamily
        Projet \projet} \\ 
      {\itshape \titre}
    \end{tabular}}
  \hfill 
  \fbox{\begin{tabular}{l}
      {~\hfill \bfseries \sffamily Groupe \groupe\ - Equipe \equipe
        \hfill~} \\[2mm] 
      Responsable : \responsible \\
      Secrétaire : \secretary \\
      Codeurs : \others
    \end{tabular}}
  \vskip 4mm ~

  ~~~\parbox{0.95\textwidth}{\small \textit{Résumé~:} \sffamily  L’objectif de ce projet est de modéliser l’écoulement de l’air autour d’un profil d’aile et d’en déduire 
  une carte de pression au-dessus et en dessous de l’aile. Cette analyse permettra d’estimer la portance générée par le profil étudié.Pour ce faire, deux étapes
   préliminaires sont nécessaires l'interpolation du profil de l'aile par la méthode des spline cubiques, puis la méthode d'intégration.}
  \vskip 1mm ~
\end{minipage}
